\documentclass[11pt,a4paper]{article}

\newcommand{\numeroRoteiro}{03}
\newcommand{\tituloCurto}{Do Dado Piloto ao Artigo no Overleaf}

% =====================================================
% PACOTES
% =====================================================
\usepackage[utf8]{inputenc}
\usepackage[T1]{fontenc}
\usepackage[brazil]{babel}
\shorthandoff{"}
\usepackage[default]{sourcesanspro}
\usepackage[scale=0.86]{sourcecodepro}
\usepackage{geometry}
\usepackage{graphicx}
\usepackage{fancyhdr}
\usepackage{titlesec}
\usepackage[table]{xcolor}
\usepackage{hyperref}
\usepackage{enumitem}
\usepackage[most]{tcolorbox}
\usepackage{float}
\usepackage{tikz}
\usetikzlibrary{arrows.meta, positioning, shapes.geometric, calc, babel}
\usepackage{tabularx}
\usepackage{booktabs}
\usepackage{listings}
\usepackage{lastpage}
\usepackage{qrcode}
\usepackage{multicol}
\usepackage{amssymb}

\geometry{top=2.6cm, bottom=2.6cm, left=2cm, right=2cm, headsep=0.8cm}
\linespread{1.18}
\setlength{\parskip}{0.55em}
\setlength{\parindent}{0pt}

% =====================================================
% PALETA (cores institucionais IFSC, em tons suaves)
% =====================================================
\definecolor{azul}{RGB}{30,60,90}
\definecolor{azulMedio}{RGB}{37,99,235}
\definecolor{verde}{RGB}{22,101,52}
\definecolor{verdeVivo}{RGB}{22,163,74}
\definecolor{ambar}{RGB}{217,119,6}
\definecolor{vermelho}{RGB}{220,38,38}
\definecolor{roxo}{RGB}{124,58,237}
\definecolor{tinta}{RGB}{30,41,59}
\definecolor{suave}{RGB}{100,116,139}
\definecolor{borda}{RGB}{226,232,240}
\definecolor{fundo}{RGB}{248,250,252}
\definecolor{barraTitulo}{RGB}{241,245,249}
\definecolor{terminal}{RGB}{15,23,42}

\color{tinta}

% =====================================================
% ÍCONES VETORIAIS TIKZ (Compatíveis com Tectonic e PDFLaTeX)
% =====================================================
\newcommand{\icoTable}{\tikz[scale=0.08, baseline=-2pt]{\draw[line width=0.8pt] (0,0) grid (3,2);}}
\newcommand{\icoPlay}{\tikz[scale=0.10, baseline=-2pt]{\fill (0,0) -- (0,1.8) -- (1.6,0.9) -- cycle;}}
\newcommand{\icoChart}{\tikz[scale=0.08, baseline=-2pt]{\fill (0,0) rectangle (0.8,1.2); \fill (1.1,0) rectangle (1.9,2.2); \fill (2.2,0) rectangle (3.0,1.6);}}
\newcommand{\icoUpload}{\tikz[scale=0.10, baseline=-2pt]{\draw[line width=0.9pt, -{Latex[length=1.5mm]}] (1,0) -- (1,1.8); \draw[line width=0.8pt] (0,0) -- (2,0);}}
\newcommand{\icoCheck}{\tikz[scale=0.10, baseline=-2pt]{\draw[line width=1.1pt, line cap=round] (0,0.8) -- (0.7,0.1) -- (2,1.6);}}
\newcommand{\icoCheckDouble}{\tikz[scale=0.09, baseline=-2pt]{\draw[line width=1pt] (0,0.8) -- (0.6,0.1) -- (1.6,1.4); \draw[line width=1pt] (0.6,0.8) -- (1.2,0.1) -- (2.2,1.4);}}
\newcommand{\icoBulb}{\tikz[scale=0.09, baseline=-2pt]{\draw[line width=0.8pt] (1,1.8) circle (0.8); \fill (0.6,0.7) rectangle (1.4,0.4);}}
\newcommand{\icoWarn}{\tikz[scale=0.10, baseline=-2pt]{\draw[line width=0.9pt] (1,2) -- (0,0) -- (2,0) -- cycle; \node at (1,0.8) {\bfseries\tiny !};}}
\newcommand{\icoBook}{\tikz[scale=0.09, baseline=-2pt]{\draw[line width=0.8pt] (0,0) rectangle (1.6,2); \draw[line width=0.7pt] (0.8,0) -- (0.8,2);}}
\newcommand{\icoUsers}{\tikz[scale=0.08, baseline=-2pt]{\draw[line width=0.8pt] (1,2) circle (0.6); \draw[line width=0.8pt] (0.2,0) arc(180:0:0.8); \draw[line width=0.8pt] (2.4,1.8) circle (0.5); \draw[line width=0.8pt] (1.8,0) arc(180:0:0.6);}}
\newcommand{\icoTarget}{\tikz[scale=0.08, baseline=-2pt]{\draw[line width=0.8pt] (1,1) circle (1.2); \draw[line width=0.8pt] (1,1) circle (0.6); \fill (1,1) circle (0.2);}}
\newcommand{\icoDoc}{\tikz[scale=0.09, baseline=-2pt]{\draw[line width=0.8pt] (0,0) -- (0,2) -- (1.2,2) -- (1.8,1.4) -- (1.8,0) -- cycle; \draw[line width=0.7pt] (1.2,2) -- (1.2,1.4) -- (1.8,1.4);}}
\newcommand{\icoGrad}{\tikz[scale=0.09, baseline=-2pt]{\fill (0,1.2) -- (1.5,2) -- (3,1.2) -- (1.5,0.4) -- cycle; \draw[line width=0.8pt] (0.7,0.8) -- (0.7,0.2) arc(180:360:0.8) -- (2.3,0.8);}}
\newcommand{\icoTeacher}{\tikz[scale=0.09, baseline=-2pt]{\draw[line width=0.8pt] (0,0) rectangle (2.2,1.6); \draw[line width=0.8pt] (1.1,0) -- (1.1,-0.6); \draw[line width=0.8pt] (0.5,-0.6) -- (1.7,-0.6);}}
\newcommand{\icoSquare}{\tikz[scale=0.10, baseline=-1pt]{\draw[line width=0.8pt] (0,0) rectangle (1.4,1.4);}}
\newcommand{\icoChevron}{\tikz[scale=0.08, baseline=-1pt]{\draw[line width=0.9pt] (0,1.6) -- (1,0.8) -- (0,0);}}
\newcommand{\icoStream}{\tikz[scale=0.08, baseline=-2pt]{\draw[line width=1pt] (0,2) -- (2,2); \draw[line width=1pt] (0,1) -- (2,1); \draw[line width=1pt] (0,0) -- (2,0);}}

% =====================================================
% LISTINGS
% =====================================================
\lstdefinestyle{base}{
    basicstyle=\ttfamily\small\color{tinta},
    keywordstyle=\color{azulMedio}\bfseries,
    commentstyle=\color{suave}\itshape,
    stringstyle=\color{verde},
    breaklines=true, columns=fullflexible, keepspaces=true,
    showstringspaces=false, upquote=true, tabsize=4,
    frame=none, numbers=none, aboveskip=0pt, belowskip=0pt,
    extendedchars=true,
    literate=
      {á}{{\'a}}1 {é}{{\'e}}1 {í}{{\'i}}1 {ó}{{\'o}}1 {ú}{{\'u}}1
      {Á}{{\'A}}1 {É}{{\'E}}1 {Í}{{\'I}}1 {Ó}{{\'O}}1 {Ú}{{\'U}}1
      {ã}{{\~a}}1 {õ}{{\~o}}1 {Ã}{{\~A}}1 {Õ}{{\~O}}1
      {â}{{\^a}}1 {ê}{{\^e}}1 {ô}{{\^o}}1 {Â}{{\^A}}1 {Ê}{{\^E}}1 {Ô}{{\^O}}1
      {à}{{\`a}}1 {ç}{{\c c}}1 {Ç}{{\c C}}1
      {–}{{--}}1 {±}{{$\pm$}}1 {≈}{{$\approx$}}1 {÷}{{$\div$}}1
}
\lstdefinestyle{py}{style=base, language=Python, morekeywords={True,False,None}}
\lstdefinestyle{tex}{style=base, language=TeX}
\lstdefinestyle{txt}{style=base}
\lstset{style=base}

% =====================================================
% COMPONENTES VISUAIS
% =====================================================
\newcommand{\pontosJanela}{%
  {\color{vermelho!70}\tikz\fill (0,0) circle (2.2pt);}\,%
  {\color{ambar!80}\tikz\fill (0,0) circle (2.2pt);}\,%
  {\color{verdeVivo!70}\tikz\fill (0,0) circle (2.2pt);}\hspace{0.8em}}

\newtcblisting{codigo}[2][py]{
    enhanced, listing only, breakable,
    colback=fundo, colframe=borda, boxrule=0.6pt, arc=6pt,
    left=10pt, right=10pt, top=6pt, bottom=6pt,
    title={\pontosJanela #2}, coltitle=suave, colbacktitle=barraTitulo,
    fonttitle=\ttfamily\footnotesize, toptitle=2pt, bottomtitle=2pt,
    titlerule=0.6pt, titlerule style=borda,
    listing options={style=#1}
}

\newtcolorbox{destaque}[3]{
    enhanced, breakable,
    colback=#1!6, frame hidden, arc=6pt,
    borderline west={3.5pt}{0pt}{#1},
    left=14pt, right=12pt, top=9pt, bottom=9pt,
    before upper={{\color{#1}\large #2}\hspace{0.55em}{\color{#1!75!black}\bfseries #3}\par\vspace{0.25em}}
}
\newenvironment{dica}[1]{\begin{destaque}{ambar}{\icoBulb}{#1}}{\end{destaque}}
\newenvironment{atencao}[1]{\begin{destaque}{vermelho}{\icoWarn}{#1}}{\end{destaque}}
\newenvironment{conceito}[1]{\begin{destaque}{azulMedio}{\icoBook}{#1}}{\end{destaque}}
\newenvironment{sucesso}[1]{\begin{destaque}{verdeVivo}{\icoCheck}{#1}}{\end{destaque}}

\newtcolorbox{painel}[3][]{
    enhanced, breakable,
    colback=white, colframe=borda, boxrule=0.6pt, arc=8pt,
    left=14pt, right=14pt, top=14pt, bottom=10pt,
    title={#2\hspace{0.5em}#3}, coltitle=white, colbacktitle=azul,
    fonttitle=\bfseries, attach boxed title to top left={xshift=12pt, yshift=-8pt},
    boxed title style={arc=4pt, boxrule=0pt}, #1
}

\tcbset{cartao/.style={enhanced, colback=fundo, colframe=borda, boxrule=0.6pt, arc=8pt,
    left=10pt, right=10pt, top=8pt, bottom=8pt, fontupper=\small\raggedright}}

\newcommand{\tecla}[1]{\tcbox[on line, enhanced, colback=white, colframe=suave!50, boxrule=0.5pt,
    arc=2.5pt, left=3pt, right=3pt, top=0.5pt, bottom=0.5pt, boxsep=0.5pt,
    borderline south={1.2pt}{0pt}{suave!40}]{\ttfamily\footnotesize #1}}
\newcommand{\botao}[1]{\tcbox[on line, enhanced, colback=azulMedio!8, colframe=azulMedio!35, boxrule=0.5pt,
    arc=3pt, left=3pt, right=3pt, top=0.5pt, bottom=0.5pt, boxsep=0.5pt]{\sffamily\footnotesize\color{azul}#1}}
\newcommand{\seta}{\,{\color{suave}\scriptsize\icoChevron}\,}
\newcommand{\chip}[2]{\tcbox[on line, enhanced, colback=white, opacityback=0.14, colframe=white,
    opacityframe=0.45, boxrule=0.5pt, arc=8pt, left=5pt, right=5pt, top=1.5pt, bottom=1.5pt,
    boxsep=0pt, coltext=white]{\footnotesize #1\hspace{0.4em}#2}}
\newcommand{\circulo}[2][azul]{\tikz[baseline=(c.base)]\node[circle, fill=#1, text=white,
    inner sep=0pt, minimum size=1.35em, font=\footnotesize\bfseries](c){#2};}
\newcommand{\caixaMarcar}{{\color{suave}\icoSquare}}

\hypersetup{colorlinks=true, linkcolor=azulMedio, urlcolor=verde, citecolor=azul}
\newcommand{\urlDashboard}{https://chameoandre.github.io/CST-Sistemas-para-internet-iniciacao-cientifica-extensionista/2026-2/}
\newcommand{\urlColab}{https://colab.research.google.com/github/chameoandre/CST-Sistemas-para-internet-iniciacao-cientifica-extensionista/blob/main/2026-2/roteiro-2-python-graficos/analise_experimento.ipynb}

% =====================================================
% CABEÇALHO E RODAPÉ
% =====================================================
\pagestyle{fancy}
\fancyhf{}
\fancyhead[L]{%
    \IfFileExists{figures/garopaba_horizontal_marca2015_PNG.png}{%
        \raisebox{-0.25\height}{\includegraphics[height=0.95cm]{figures/garopaba_horizontal_marca2015_PNG.png}}%
    }{\textbf{\color{verde} IFSC — Câmpus Garopaba}}}
\fancyhead[R]{\footnotesize\color{suave}\textbf{\color{azul}Roteiro Prático 03}\hspace{0.5em}\textbullet\hspace{0.5em}\tituloCurto}
\fancyfoot[L]{\footnotesize\color{suave}Iniciação Científica e Extensionista\hspace{0.4em}\textbullet\hspace{0.4em}CST Sistemas para a Internet}
\fancyfoot[R]{\tcbox[on line, colback=azul!8, colframe=azul!8, arc=6pt, boxrule=0pt, left=5pt, right=5pt, top=1pt, bottom=1pt, boxsep=0pt]{\footnotesize\color{azul}\bfseries\thepage\,/\,\pageref*{LastPage}}}
\renewcommand{\headrule}{{\color{borda}\hrule height 0.6pt}}
\renewcommand{\footrulewidth}{0pt}
\setlength{\headheight}{36pt}

% =====================================================
% TÍTULOS DE SEÇÃO
% =====================================================
\titleformat{\section}
  {\normalfont\Large\bfseries\color{azul}}
  {\tikz[baseline=(n.base)]\node[fill=azul, rounded corners=5pt, text=white, inner sep=5pt,
      minimum width=1.9em, font=\large\bfseries](n){\thesection};}
  {0.7em}{}
\titlespacing*{\section}{0pt}{1.5em}{0.6em}

\titleformat{\subsection}
  {\normalfont\large\bfseries\color{azul!85!black}}
  {{\color{verdeVivo}\thesubsection}}
  {0.6em}{}
\titlespacing*{\subsection}{0pt}{1.1em}{0.4em}

\setlist[itemize]{leftmargin=1.4em, itemsep=0.25em, topsep=0.3em, label={\color{verdeVivo}\small\icoChevron}}
\setlist[enumerate]{leftmargin=2.2em, itemsep=0.45em, topsep=0.3em}

\newcolumntype{Y}{>{\raggedright\arraybackslash}X}
\newcommand{\cabecalho}{\rowcolor{azul}\color{white}\bfseries}
\newcommand{\cab}[1]{\textcolor{white}{\textbf{#1}}}

\begin{document}

% =====================================================
% FAIXA DE ABERTURA (CAPA)
% =====================================================
\begin{tcolorbox}[enhanced, breakable=false,
    colback=azul, colframe=azul, arc=12pt,
    interior style={left color=azul, right color=verde},
    left=18pt, right=18pt, top=16pt, bottom=16pt]
    {\color{white!80!black}\footnotesize\bfseries\icoDoc\hspace{0.5em}ROTEIRO PRÁTICO 03 \hspace{0.4em}\textbullet\hspace{0.4em} OFICINA DE INTEGRAÇÃO EXPERIMENTAL}\par\vspace{3pt}
    {\color{white}\fontsize{21pt}{25pt}\selectfont\bfseries Do Dado Piloto ao Artigo no Overleaf}\par\vspace{3pt}
    {\color{white!90!black}\small Pipeline Completo: Planilha Universal, Starter Kit no Colab e Inclusão no Template COTB 2027}\par\vspace{10pt}
    \chip{\icoUsers}{Duplas ou Individual}\hspace{0.4em}
    \chip{\icoChart}{Google Colab + Python}\hspace{0.4em}
    \chip{\icoDoc}{LaTeX Overleaf (ACM)}
\end{tcolorbox}

\vspace{0.4em}

% =====================================================
% TRÊS CARTÕES INSTITUCIONAIS
% =====================================================
\begin{tcbraster}[raster columns=3, raster equal height, raster column skip=8pt]
\begin{tcolorbox}[cartao]
{\color{azulMedio}\icoGrad}\hspace{0.4em}\textbf{Unidade Curricular}\par\vspace{2pt}
Iniciação Científica e Extensionista (2ª Fase / 60h) \textbullet\ CST Sistemas para a Internet.
\end{tcolorbox}
\begin{tcolorbox}[cartao]
{\color{verdeVivo}\icoTeacher}\hspace{0.4em}\textbf{Equipe Docente}\par\vspace{2pt}
Prof. André Moraes \& Prof. Adalberto Tabalipa (Mediação em sala: Prof. Alberto).
\end{tcolorbox}
\begin{tcolorbox}[cartao]
{\color{ambar}\icoTarget}\hspace{0.4em}\textbf{Meta do Encontro}\par\vspace{2pt}
Sair da aula com a planilha da dupla rodada e os gráficos compilados dentro do artigo no Overleaf.
\end{tcolorbox}
\end{tcbraster}

% =====================================================
% PAINEL DE OBJETIVOS
% =====================================================
\begin{painel}{\icoTarget}{Ao final desta prática você será capaz de:}
\begin{itemize}
    \item \textbf{Estruturar dados quantitativos} no formato universal de 3 colunas (\texttt{grupo}, \texttt{rodada}, \texttt{valor}) para qualquer tema de pesquisa;
    \item \textbf{Executar o Starter Kit Python} no Google Colab, adaptando as variáveis de configuração sem necessidade de programar código do zero;
    \item \textbf{Exportar gráficos científicos em PDF} no tamanho exato de uma coluna ACM (3,33 polegadas) e a tabela em formato \texttt{booktabs};
    \item \textbf{Inserir as figuras e a tabela no Overleaf} (\texttt{\textbackslash begin\{figure\}} e \texttt{\textbackslash input\{tabela\_resultados\}}), compilando o artigo sem erros;
    \item \textbf{Redigir a interpretação dos resultados} com foco na pergunta da pesquisa e no impacto social da extensão.
\end{itemize}
\end{painel}

% =====================================================
% DIAGRAMA DE ANDAMENTO NUMERADO COM ÍCONES
% =====================================================
\section{Diagrama de Andamento da Prática}\label{sec:diagrama}

A Figura~\ref{fig:fluxo_andamento} sintetiza o fluxo de trabalho contínuo do laboratório. Observe a numeração sequencial de 1 a 5 e os ícones informativos de cada etapa:

\begin{figure}[H]
\centering
\begin{tikzpicture}[
    etapaNode/.style={
        rectangle, draw=borda, fill=white, line width=0.8pt,
        rounded corners=8pt, text width=2.65cm, minimum height=3.0cm,
        align=center, font=\footnotesize
    },
    numNode/.style={
        circle, fill=azul, text=white, inner sep=0pt, minimum size=1.55em,
        font=\footnotesize\bfseries
    },
    setaFluxo/.style={
        draw=suave!80, line width=1.3pt, -{Latex[length=3.0mm, width=2.2mm]}
    }
]

% Nó 1: Planilha Universal
\node[etapaNode, fill=fundo] (n1) {
    \vspace{4pt}
    {\color{azulMedio}\fontsize{18pt}{18pt}\selectfont\icoTable}\\[6pt]
    \textbf{1. Planilha Piloto}\\[2pt]
    {\color{suave}\scriptsize 3 colunas padrão:\\[1pt]
    \texttt{grupo, rodada, valor}\\[2pt]
    Salvar em \texttt{.csv}}
};
\node[numNode, fill=azulMedio] at ($(n1.north)+(0,0.1)$) {1};

% Nó 2: Colab
\node[etapaNode, right=0.45cm of n1, fill=verdeVivo!5, draw=verdeVivo!40] (n2) {
    \vspace{4pt}
    {\color{verdeVivo}\fontsize{18pt}{18pt}\selectfont\icoPlay}\\[6pt]
    \textbf{2. Google Colab}\\[2pt]
    {\color{suave}\scriptsize Subir o CSV na pasta\\[1pt]
    Trocar 3 variáveis\\[2pt]
    Executar Tudo}
};
\node[numNode, fill=verdeVivo] at ($(n2.north)+(0,0.1)$) {2};

% Nó 3: Artefatos
\node[etapaNode, right=0.45cm of n2, fill=fundo] (n3) {
    \vspace{4pt}
    {\color{ambar}\fontsize{18pt}{18pt}\selectfont\icoChart}\\[6pt]
    \textbf{3. Artefatos}\\[2pt]
    {\color{suave}\scriptsize Baixar do Colab:\\[1pt]
    \texttt{grafico\_barras.pdf}\\[1pt]
    \texttt{tabela\_res.tex}}
};
\node[numNode, fill=ambar] at ($(n3.north)+(0,0.1)$) {3};

% Nó 4: Overleaf
\node[etapaNode, right=0.45cm of n3, fill=roxo!5, draw=roxo!35] (n4) {
    \vspace{4pt}
    {\color{roxo}\fontsize{18pt}{18pt}\selectfont\icoUpload}\\[6pt]
    \textbf{4. Overleaf}\\[2pt]
    {\color{suave}\scriptsize Upload no projeto COTB\\[1pt]
    Inserir \texttt{\textbackslash includegraphics}\\[1pt]
    e \texttt{\textbackslash input}}
};
\node[numNode, fill=roxo] at ($(n4.north)+(0,0.1)$) {4};

% Nó 5: Compilação & Registro
\node[etapaNode, right=0.45cm of n4, fill=verdeVivo!7, draw=verdeVivo] (n5) {
    \vspace{4pt}
    {\color{verde}\fontsize{18pt}{18pt}\selectfont\icoCheckDouble}\\[6pt]
    \textbf{5. Artigo Pronto}\\[2pt]
    {\color{suave}\scriptsize Clicar \textit{Recompile}\\[1pt]
    Conferir artigo ACM\\[1pt]
    Link no Dashboard}
};
\node[numNode, fill=verde] at ($(n5.north)+(0,0.1)$) {5};

% Conexões com setas
\draw[setaFluxo] (n1) -- (n2);
\draw[setaFluxo] (n2) -- (n3);
\draw[setaFluxo] (n3) -- (n4);
\draw[setaFluxo] (n4) -- (n5);

\end{tikzpicture}
\caption{Diagrama de andamento numerado da oficina de integração experimental.}
\label{fig:fluxo_andamento}
\end{figure}

% =====================================================
% ETAPA 1: PLANILHA UNIVERSAL
% =====================================================
\section{\circulo{1}\hspace{0.4em}Etapa 1: A Planilha Universal de 3 Colunas}\label{sec:etapa1}

Para que o script funcione imediatamente sem erros de sintaxe ou colunas não encontradas, todas as duplas seguirão o \textbf{formato universal de três colunas}:

\begin{codigo}[txt]{modelo\_dados\_piloto.csv (estrutura obrigatória)}
grupo,rodada,valor
Grupo A,1,85.5
Grupo A,2,82.0
...
Grupo B,1,62.0
...
Grupo C,5,93.5
\end{codigo}

\begin{conceito}{Como adaptar o seu tema para a planilha?}
\begin{itemize}
    \item \textbf{Coluna \texttt{grupo}:} Representa as tecnologias, ferramentas ou categorias comparadas na sua pesquisa (ex.: \texttt{Moodle}, \texttt{Canvas}, \texttt{Classroom} para AVAs; ou \texttt{WebRTC}, \texttt{HLS}, \texttt{SRT} para Telemedicina; ou \texttt{Sinal A}, \texttt{Sinal B} para Libras).
    \item \textbf{Coluna \texttt{rodada}:} O número do teste ou medição repetida (1 a 5).
    \item \textbf{Coluna \texttt{valor}:} Apenas números com ponto decimal (nota WCAG, taxa de acurácia, latência em ms ou tempo em segundos).
\end{itemize}
\end{conceito}

\begin{dica}{Exportação correta do Google Planilhas}
Ao terminar de preencher no Google Planilhas, acesse o menu superior:\par\vspace{2pt}
\botao{Arquivo} \seta \botao{Fazer download} \seta \botao{Valores separados por vírgula (.csv)}.\par\vspace{2pt}
Mantenha a extensão \texttt{.csv} e não converta para PDF ou Excel.
\end{dica}

% =====================================================
% ETAPA 2: GOOGLE COLAB
% =====================================================
\section{\circulo{2}\hspace{0.4em}Etapa 2: Execução no Starter Kit do Google Colab}\label{sec:etapa2}

Abra o notebook oficial da disciplina através do link direto do Dashboard ou do GitHub:
\begin{center}
\href{\urlColab}{\botao{\icoPlay\hspace{0.4em}\textbf{Clique aqui para abrir o Starter Kit no Google Colab}}}
\end{center}

\begin{enumerate}
    \item No painel esquerdo do Colab, clique no ícone da pasta e faça o upload do seu arquivo \texttt{.csv};
    \item Localize a célula de \textbf{CONFIGURAÇÃO} e altere apenas o nome do arquivo e a unidade:
\end{enumerate}

\begin{codigo}{Ajuste de apenas 4 linhas no Google Colab}
ARQUIVO_DADOS  = "meus_dados.csv"   # Nome exato do seu arquivo enviado
COLUNA_GRUPO   = "grupo"            # Mantém "grupo"
COLUNA_METRICA = "valor"            # Mantém "valor"
ROTULO_EIXO_Y  = "Acurácia Média (%)"  # Nome da sua métrica + unidade!
\end{codigo}

\begin{enumerate}[resume]
    \item No menu superior do Colab, clique em \botao{Ambiente de execução}\seta\botao{Executar tudo} (ou use o atalho \tecla{Ctrl} + \tecla{F9}).
\end{enumerate}

% =====================================================
% ETAPA 3: EXPORTAÇÃO DE ARTEFATOS
% =====================================================
\section{\circulo{3}\hspace{0.4em}Etapa 3: Download dos Artefatos Gerados}\label{sec:etapa3}

Após a execução completa, observe que a pasta lateral contém três novos arquivos de publicação científica:
\begin{itemize}
    \item \texttt{grafico\_barras.pdf}: Gráfico vetorial de barras com Média $\pm$ Desvio Padrão na largura exata de 1 coluna ACM (3,33 polegadas);
    \item \texttt{grafico\_boxplot.pdf}: Gráfico de caixa demonstrando a dispersão das 5 rodadas;
    \item \texttt{tabela\_resultados.tex}: Tabela em código \LaTeX{} com o pacote \texttt{booktabs} pronta para o artigo.
\end{itemize}

Clique com o botão direito sobre cada um deles e selecione \botao{Fazer download}.

% =====================================================
% ETAPA 4: ENXERTO NO OVERLEAF
% =====================================================
\section{\circulo{4}\hspace{0.4em}Etapa 4: Inclusão no Artigo no Overleaf}\label{sec:etapa4}

Abra o projeto do artigo da dupla no \textbf{Overleaf} (template oficial ACM \texttt{sigconf} do COTB 2027):

\begin{enumerate}
    \item No menu lateral esquerdo do Overleaf, clique em \botao{Upload} e envie o arquivo \texttt{grafico\_barras.pdf} e o arquivo \texttt{tabela\_resultados.tex};
    \item Abra o arquivo principal \texttt{main.tex} e localize a seção de \textbf{Resultados};
    \item Insira o bloco de código a seguir:
\end{enumerate}

\begin{codigo}[tex]{main.tex (seção de Resultados)}
\section{Resultados e Discussão}

A Figura~\ref{fig:desempenho} apresenta as médias observadas nos testes
experimentais, enquanto a Tabela~\ref{tab:resultados} detalha as
estatísticas descritivas das rodadas de medição.

\begin{figure}[htbp]
  \centering
  \includegraphics[width=\columnwidth]{grafico_barras.pdf}
  \caption{Média e desvio padrão das medições de desempenho por grupo testado (n = 5 rodadas).}
  \label{fig:desempenho}
\end{figure}

\input{tabela_resultados}
\end{codigo}

\begin{atencao}{Regras de ouro no LaTeX}
\begin{itemize}
    \item O comando \texttt{\textbackslash label\{fig:...\}} deve ficar \textbf{sempre depois} do \texttt{\textbackslash caption\{...\}};
    \item Use sempre \texttt{width=\textbackslash columnwidth} para que a figura ocupe exatamente a largura de uma coluna sem distorção;
    \item Não digite a tabela à mão: use \texttt{\textbackslash input\{tabela\_resultados\}} para importar o arquivo gerado pelo Python.
\end{itemize}
\end{atencao}

% =====================================================
% ETAPA 5: COMPILAÇÃO E REGISTRO
% =====================================================
\section{\circulo{5}\hspace{0.4em}Etapa 5: Recompilação \& Registro no Painel}\label{sec:etapa5}

\begin{enumerate}
    \item Clique no botão verde \botao{Recompile} no Overleaf;
    \item Verifique se o artigo compilou sem mensagens de erro e se a figura e a tabela aparecem perfeitamente alinhadas no PDF do artigo;
    \item No canto superior direito do Overleaf, clique em \botao{Share} e copie o link de leitura (\texttt{/read/...});
    \item Acesse o \textbf{Dashboard da UC} (\href{\urlDashboard}{\textbf{clique aqui para abrir}}), vá na \textbf{Aba 4 (Painel de Equipes)}, localize o card da sua dupla e atualize o link do Overleaf e o status para o \textbf{Checkpoint 2}.
\end{enumerate}

% =====================================================
% CHECKLIST DE VERIFICAÇÃO
% =====================================================
\vspace{0.8em}
\begin{tcolorbox}[enhanced, breakable, colback=verdeVivo!5, colframe=verdeVivo!35, boxrule=0.6pt, arc=8pt, left=14pt, right=14pt, top=10pt, bottom=10pt]
{\color{verde}\bfseries\large\icoCheck\hspace{0.5em}Checklist de Conclusão da Oficina}\par\vspace{4pt}
\begin{itemize}[label={\caixaMarcar}, itemsep=0.4em]
    \item Planilha criada com 3 colunas obrigatórias (\texttt{grupo}, \texttt{rodada}, \texttt{valor}) e exportada em \texttt{.csv};
    \item Notebook Colab executado sem erros com os dados da dupla;
    \item \texttt{grafico\_barras.pdf} e \texttt{tabela\_resultados.tex} baixados com sucesso;
    \item Arquivos enviados ao Overleaf e referenciados no \texttt{main.tex};
    \item Artigo compilando perfeitamente em 2 colunas com a figura e tabela visíveis;
    \item Link de leitura do Overleaf registrado no Painel de Equipes do Dashboard.
\end{itemize}
\end{tcolorbox}

\end{document}
